% !TeX root = ../main.tex
% Fonte: Lezione 13, pp. 5--13. Casi costanti, quozienti generali e svolgimenti completati.
\section{Limiti di polinomi e termine dominante}
\label{sec:polinomi-limiti}

\begin{proposizione}
Per ogni $k\in\N^*$ si ha $n^k\to+\infty$ e, di conseguenza, $1/n^k\to0^+$.
\end{proposizione}
\begin{proof}
Si procede per induzione su $k$. Per $k=1$, $n\to+\infty$ è già stato verificato dalla definizione. Se $n^k\to+\infty$, allora $n^{k+1}=n^k\cdot n\to+\infty$ per la regola del prodotto. Il limite dei reciproci segue dalla proposizione~\ref{prop:limite-reciproco}.
\end{proof}

\begin{esempio}
La successione $3n^3-n+2$ presenta, se si considerano separatamente gli addendi, la forma indeterminata $+\infty-\infty$. Non si può dedurre il risultato dai soli limiti degli addendi. Il termine $3n^3$ cresce però più rapidamente di $n$: raccoglierlo permette di confrontare le due grandezze.
\end{esempio}

\begin{proposizione}[Termine dominante di un polinomio]
\label{prop:polinomio-dominante}
Sia $P(x)=c_0+c_1x+\cdots+c_kx^k$ un polinomio di grado $k\ge1$, con $c_k\ne0$. Allora
\[
\frac{P(n)}{c_kn^k}\longrightarrow1,
\qquad
P(n)\longrightarrow
\begin{cases}
+\infty,&c_k>0,\\
-\infty,&c_k<0.
\end{cases}
\]
Un polinomio costante ha invece come limite il proprio valore, compreso il polinomio nullo.
\end{proposizione}
\begin{proof}
\textbf{Raccoglimento.} Per $n\ge1$,
\[
P(n)=c_kn^k\left(1+\sum_{j=0}^{k-1}\frac{c_j}{c_kn^{k-j}}\right).
\]
Per ogni $j\le k-1$ l'esponente $k-j$ è un intero positivo, dunque $1/n^{k-j}\to0$. La somma finita nella parentesi tende a $0$, per l'algebra dei limiti. Si ottiene così $P(n)/(c_kn^k)\to1$.

\textbf{Segno del limite.} La parentesi tende a $1$, mentre $c_kn^k$ tende a $+\infty$ se $c_k>0$ e a $-\infty$ se $c_k<0$. La regola del prodotto dà la conclusione.
\end{proof}

Se i coefficienti non nulli dei termini di grado positivo hanno tutti lo stesso segno, si può usare direttamente la regola della somma. Il raccoglimento funziona anche quando i segni sono diversi: è il coefficiente di grado massimo a decidere il limite.

\begin{esempio}
Per il polinomio precedente,
\[
3n^3-n+2=3n^3\left(1-\frac1{3n^2}+\frac2{3n^3}\right)\longrightarrow+\infty.
\]
Analogamente,
\[
3n^3-2n^2+4n+1
=3n^3\left(1-\frac2{3n}+\frac4{3n^2}+\frac1{3n^3}\right)
\longrightarrow+\infty.
\]
In entrambi i casi la parentesi tende a $1$; non si sta assegnando un valore alla forma $+\infty-\infty$.
\end{esempio}

\subsection{La notazione o piccolo di uno}
% Fonte: Lezione 13, pp. 11--12; esplicitato il significato della notazione usata nella fonte.

La scrittura $r_n=o(1)$ significa semplicemente che $r_n\to0$. Il risultato sul termine dominante si abbrevia quindi in
\[
P(n)=c_kn^k\bigl(1+o(1)\bigr).
\]
Il simbolo $o(1)$ rappresenta una successione infinitesima; due sue occorrenze possono indicare successioni diverse. Non è un numero da cancellare fra numeratore e denominatore: si usano i limiti delle espressioni che lo contengono.

\Needspace{22\baselineskip}
\section{Quozienti di polinomi}
% Fonte: Lezione 13, pp. 11--12; regola generale ricostruita dai raccoglimenti.

\begin{proposizione}[Confronto dei gradi]
Siano $P$ e $Q$ polinomi non nulli di gradi $k,h\ge0$ e coefficienti principali $c_k,d_h$. Il quoziente $P(n)/Q(n)$ è definito definitivamente e
\[
\frac{P(n)}{Q(n)}\longrightarrow
\begin{cases}
0,&k<h,\\
c_k/d_h,&k=h,\\
+\infty,&k>h\text{ e }c_k/d_h>0,\\
-\infty,&k>h\text{ e }c_k/d_h<0.
\end{cases}
\]
Se $P$ è il polinomio nullo e $Q$ non lo è, il quoziente è definitivamente nullo.
\end{proposizione}
\begin{proof}
Il raccoglimento dà $P(n)=c_kn^k(1+r_n)$ e $Q(n)=d_hn^h(1+s_n)$, con $r_n,s_n\to0$; per un polinomio costante si pone il resto uguale a zero. Poiché $1+s_n\to1$, definitivamente $1+s_n>1/2$, dunque $Q(n)\ne0$. Si ha
\[
\frac{P(n)}{Q(n)}
=\frac{c_k}{d_h}\,n^{k-h}\,\frac{1+r_n}{1+s_n},
\qquad \frac{1+r_n}{1+s_n}\longrightarrow1.
\]
Se $k<h$, il fattore $n^{k-h}=1/n^{h-k}$ tende a $0$; se $k=h$ è identicamente $1$; se $k>h$ tende a $+\infty$. La regola del prodotto conclude, tenendo conto del segno di $c_k/d_h$.
\end{proof}

\begin{esempio}
Per $P(n)=3n^3-n+1$ e $Q(n)=2n^4+n+1$, entrambi i polinomi tendono a $+\infty$: la regola diretta del quoziente non risolve la forma $+\infty/+\infty$. Raccogliendo separatamente,
\[
\frac{3n^3-n+1}{2n^4+n+1}
=\frac3{2n}\,
\frac{1-\frac1{3n^2}+\frac1{3n^3}}
{1+\frac1{2n^3}+\frac1{2n^4}}
\longrightarrow\frac32\cdot0\cdot1=0.
\]
Equivalentemente, il quoziente si scrive $\frac3{2n}\frac{1+o(1)}{1+o(1)}$. I due resti non sono uguali, ma il loro contributo al rapporto tende a $1$.
\end{esempio}

\begin{esercizio}[Individuare il grado effettivo]
Calcolare i limiti, per $n\to+\infty$, delle tre successioni
\[
\frac{-n^6+2n^4+3n}{-n^5+5n^4+9n^2},\qquad
\frac{-n^6+2n^4+3n}{-n^5+5n^6+9n^2},\qquad
\frac{n^{37}-5n^{51}+12}{n^{11}-7n^3+8n^{51}-9}.
\]
\end{esercizio}
\begin{svolgimento}
\textbf{Denominatore di grado cinque.} Nel primo rapporto i termini dominanti sono $-n^6$ e $-n^5$. Infatti
\[
\frac{-n^6+2n^4+3n}{-n^5+5n^4+9n^2}
=n\,\frac{1-2/n^2-3/n^5}{1-5/n-9/n^3}
\longrightarrow+\infty.
\]

\textbf{Denominatore di grado sei.} Nel secondo rapporto il termine $5n^6$ domina anche se compare dopo $-n^5$. Dividendo per $n^6$,
\[
\frac{-1+2/n^2+3/n^5}{-1/n+5+9/n^4}
\longrightarrow-\frac15.
\]
Conta il grado massimo, non l'ordine in cui sono scritti i monomi.\footnote{Le due versioni distinguono un esponente di lettura incerta nel manoscritto: il risultato cambia se il termine del denominatore è $5n^4$ oppure $5n^6$. Il raccoglimento con $-n^5$ è valido soltanto nella prima versione.}

\Needspace{9\baselineskip}
\textbf{Gradi uguali.} Nell'ultimo rapporto entrambi i gradi sono $51$. Raccogliendo i rispettivi termini principali,
\[
\frac{n^{37}-5n^{51}+12}{n^{11}-7n^3+8n^{51}-9}
=\frac{-5n^{51}(1+o(1))}{8n^{51}(1+o(1))}
\longrightarrow-\frac58.
\]
\end{svolgimento}

\section{Limite del valore assoluto}
\label{sec:limite-modulo}
% Fonte: Lezione 13, p. 13. La fonte prova il caso finito; esplicitato il caso infinito.

\begin{proposizione}
\label{prop:limite-modulo}
Se $a_n\to\ell\in\overline\R$, allora $|a_n|\to|\ell|$, dove $|+\infty|=|-\infty|=+\infty$.
\end{proposizione}
\begin{proof}
\textbf{Limite finito.} Dalla disuguaglianza triangolare si ricava
\[
|a_n|\le|a_n-\ell|+|\ell|,\qquad
|\ell|\le|a_n-\ell|+|a_n|,
\]
quindi la disuguaglianza triangolare inversa dà
\[
0\le\bigl||a_n|-|\ell|\bigr|\le|a_n-\ell|\longrightarrow0.
\]
Per i carabinieri la distanza fra $|a_n|$ e $|\ell|$ tende a zero, ossia $|a_n|\to|\ell|$.

\textbf{Limite infinito.} Se $a_n\to+\infty$, definitivamente $|a_n|=a_n$; se $a_n\to-\infty$, definitivamente $|a_n|=-a_n$. In entrambi i casi $|a_n|\to+\infty$.
\end{proof}

Per $\ell$ finito il risultato si legge anche come $f(a_n)\to f(\ell)$ per la funzione $f(x)=|x|$. Questo suggerisce di studiare quali funzioni conservino il limite delle successioni: è il punto di partenza della continuità.
